\subsection{局部常值层的积}

设 B 为拓扑空间，\((\mathcal{F}_{i})_{i\in\mathrm{I}}\) 为 B 上的一个层族。用 \(\mathcal{F}\) 表示积层 \(\prod_{i\in I}\mathcal{F}_{i}\)，并对 \(i\in I\)，设 \(pr_{i}\colon \mathcal{F}\to\mathcal{F}_{i}\) 为指标 i 的投影态射 (I, p. 46, example 7)。设 \((\mathrm{E},p)\) 与 \((\mathrm{E}_{i},p_{i})\) 为分别与层 \(\mathcal{F}\) 和 \(\mathcal{F}_{i}\) 相伴的平展 B-空间，并设 \(\varphi_{i}\colon \mathrm{E}\to \mathrm{E}_{i}\) 为与 \(pr_{i}\) 相伴的 B-态射 (I, p. 50)。最后，用 \((\mathrm{E}^{\prime},p^{\prime})\) 表示 B-积 \(\prod_{B}\mathrm{E}_{i}\)。由 B-空间积的泛性质 (I, p. 5)，存在唯一的 B-态射 \(\Phi\colon \mathrm{E}\to \mathrm{E}^{\prime}\)，使得对每个 \(i\in I\)，有 \(pr_{i}\circ\Phi=\varphi_{i}\)。

命题 10. --- 若集 I 是有限的，则 B-态射 \(\Phi\) 是同构。

由 I, p. 32 的推论，B-空间 \(\mathrm{E}'\) 是平展的，于是只须证明 B-态射 \(\Phi\) 是双射 (I, p. 30, cor. 2)。对 B 的每一点 b，限制 \(\Phi_b\colon \mathrm{E}_b\to \mathrm{E}'_b\) （\(\Phi\) 在 b 处纤维上的）与从 \(\varinjlim\prod_{i\in I}\mathcal{F}_{i}(U)\) 到 \(\prod_{i\in I}\varinjlim\mathcal{F}_{i}(U)\) 的典范映射相恒同，其中 U 取遍 B 的含 b 的开子集（参看 I, p. 50）。由 E, III, p. 67 的 prop. 10，这个映射是双射。

现在考虑拓扑空间 B 上局部常值层的情形。设 \((\mathrm{F}_{i})_{i\in\mathrm{I}}\) 为一个集族，并令 \(F=\prod_{i\in I}F_{i}\) 。如下定义典范态射 \(\psi=(\psi_{\mathrm{U}})\) （从常值层 \(\underline{F}\) 到积 \(\prod_{i\in I}\underline{F}_{i}\) （诸常值层 \(\underline{F}_{i}\) 的）中的）：对 B 的每个开子集 U 与每个局部常值函数 \(f\colon U\to F\) ，令 \(\psi_{\mathrm{U}}(f)=(\mathrm{pr}_{i}\circ f)_{i\in\mathrm{I}}\) 。

命题 11. --- 若集 I 是有限的，或空间 B 是局部连通的，则典范态射 \(\psi\colon\underline{\mathbf{F}}\to\prod\underline{\mathbf{F}_{i}}\) 是同构。

设 U 为 B 的开子集。显然映射 \(\psi_{\mathrm{U}}\) 是单射。来证明它是满射。只须证明：对每个族 \((f_i)_{i\in \mathrm{I}}\) （其中 \(f_{i}\) 是从 U 到 \(\mathrm{F}_i\) 中的局部常值映射）以及 U 的每一点 \(a\)，存在 \(a\) 在 U 中的邻域 V，使得对每个 \(i\in \mathrm{I}\)，映射 \(f_{i}|\mathrm{V}\) 都是常值的。当集 I 有限时，这样一个邻域的存在性是显然的。当空间 B 局部连通时，只须取 V 为 \(a\) 在 U 中的连通邻域。这就证明了命题。

推论 1. --- 局部常值层的有限积是局部常值的。

设 \((\mathcal{F}_{i})_{i\in\mathrm{I}}\) 为 B 上的局部常值层的有限族。B 的每一点 a 都有 B 中的开邻域 U，使得对每个 \(i\in I\)，层 \(\mathcal{F}_{i}\mid U\) 同构于常值层。于是这对层 \((\prod\mathcal{F}_{i})\mid U\) 也成立，后者等于 \(\prod(\mathcal{F}_{i}\mid U)\)。

推论 2. --- 设 \((\mathcal{F}_{i})_{i\in I}\) 为 B 上的层族。设空间 B 是局部连通的，且 B 的每一点都有开邻域 U，使得对每个 \(i\in I\) ，层 \(\mathcal{F}_{i}|U\) 同构于常值层。则积层 \(\prod \mathcal{F}_{i}\) 是局部常值的。

注. --- 1) 设 \((\mathrm{E}_{i}, p_{i})_{i \in \mathrm{I}}\) 为拓扑空间 B 的覆叠族。设空间 B 是局部连通的，且存在 B 的开覆盖 \((\mathrm{U}_{j})_{j \in \mathrm{J}}\)，使得对每个 \(j \in J\) 与每个 \(i \in I\)，覆叠 \(E_{i}\) 在 \(U_{j}\) 上均可平凡化。对 \(i \in I\)，设 \(F_{i}\) 为 \(E_{i}\) 的截面的局部常值层，设 \(F\) 为积层 \(\prod_{i} F_{i}\) 。由前一个推论，F 是局部常值层，因而与之相伴的平展 B-空间 E 是覆叠 (I, p. 86, prop. 8)。对 \(i \in I\)，设 \(pr_{i}: E \to E_{i}\) 为由指标 i 的投影态射 \(F \to F_{i}\) 诱导的 B-态射。

现在考虑 B 的一个覆叠 (Y, q)，且对每个 \(i \in I\)，设 \(f_i \colon Y \to E_i\) 为 B-态射。若用 G 表示 q 的截面的局部常值层，则诸 B-态射 \(f_i\) 诱导层态射 \(\widetilde{f}_i \colon G \to F_i\) 。设 \(\widetilde{f} \colon G \to F\) 为唯一的层态射，使得 \(\widetilde{f}_i = pr_i \circ \widetilde{f}\) 对所有 \(i \in I\) 成立。于是存在唯一的 B-态射 \(f \colon Y \to E\)，使得对每个 i 有 \(f_i = pr_i \circ f\)：它就是诱导 \(\widetilde{f}\) 的那个 B-态射。

有时称 E 为族 \((\mathrm{E}_i)_{i\in \mathrm{I}}\) 的积覆叠。

\begin{enumerate}
\def\labelenumi{\arabic{enumi})}
\setcounter{enumi}{1}
\tightlist
\item
  沿用前面的记号，典范 B-态射 \(\Phi\colon\mathrm{E}\to\Pi_{\mathrm{B}}\mathrm{E}_{i}\) 是双射。事实上该问题在 B 上是局部性的，故可设对每个 \(i\in\mathrm{I}\)，B-空间 \(\mathrm{E}_{i}\) 同构于 B-空间 \((\mathrm{B}\times\mathrm{F}_{i},\mathrm{pr}_{1})\)，其中 \(\mathrm{F}_{i}\) 是离散拓扑空间，从而层 \(\mathcal{F}_{i}\) 同构于层 \(\underline{\mathrm{F}}_{i}\)。由 prop. 11，层 \(\mathcal{F}\) 与层 \(\underline{\mathrm{F}}\) 相恒同，其中 \(\mathrm{F}\) 是集 \(\prod\mathrm{F}_{i}\)，赋以离散拓扑，且映射 \(\Phi\) 与典范映射 \(B \times F \rightarrow B \times (\prod_{i} F_{i})\) 相恒同，后者是双射。
\end{enumerate}

